% Telos shared preamble.
%
% Every Telos publication is designed to be printed on a home laser printer,
% one sheet at a time, and carried into a boat or a garage. Two constraints
% follow from that and govern everything below:
%
%   1. Pure monochrome. No value is ever a hue. Emphasis is carried by rule
%      weight, fill, and type, never by color, so a grayscale or black-only
%      printer loses nothing.
%   2. Card density. A "one-pager" is a hard contract: the leaf must fit on a
%      single side of US Letter. The layout is therefore tight by default and
%      the environments below are sized to leave room for a plate.

\documentclass[10pt,letterpaper]{article}

\usepackage[margin=0.55in,includefoot,footskip=0.24in]{geometry}
\usepackage[T1]{fontenc}
\usepackage[utf8]{inputenc}
\usepackage{lmodern}
\usepackage{microtype}
\usepackage{array}
\usepackage{booktabs}
\usepackage{longtable}
\usepackage{tabularx}
\usepackage{enumitem}
\usepackage{needspace}
\usepackage{multicol}
\usepackage{ragged2e}
\usepackage{graphicx}
\usepackage{wrapfig}
\usepackage{xcolor}
\usepackage{fancyhdr}
\usepackage{titlesec}
\usepackage{hyperref}
\usepackage[most]{tcolorbox}
\usepackage{tikz}
\usetikzlibrary{arrows.meta,calc,positioning,shapes.geometric,decorations.pathmorphing,patterns,fadings}

% Semantic names are kept so a document reads intelligibly, but every value is
% black, white, or a neutral gray that survives a black-only print driver.
\definecolor{telosink}{gray}{0}
\definecolor{telospaper}{gray}{1}
\definecolor{telosrule}{gray}{0.35}
\definecolor{telosfill}{gray}{0.90}
\definecolor{telosdeep}{gray}{0.72}
\definecolor{telosmute}{gray}{0.45}

\hypersetup{hidelinks,pdfborder={0 0 0}}

% Keep an unchanged source byte-reproducible across rebuilds: no build-clock
% creation date and no random trailer ID.
\pdfinfoomitdate=1
\pdftrailerid{}

\setlength{\parindent}{0pt}
\setlength{\parskip}{0.42em}
\setlength{\tabcolsep}{4pt}
\setlength{\columnsep}{1.1em}
\renewcommand{\arraystretch}{1.18}
\setlist[itemize]{leftmargin=1.15em,itemsep=0.12em,topsep=0.2em,parsep=0pt}
\setlist[enumerate]{leftmargin=1.5em,itemsep=0.18em,topsep=0.2em,parsep=0pt}
\setlist[description]{leftmargin=0pt,itemsep=0.16em,topsep=0.2em,parsep=0pt,
  font=\normalfont\bfseries}

% ---------------------------------------------------------------------------
% Document identity. Each leaf declares these before \begin{document} so the
% running head, the colophon, and the PDF metadata all agree.
% ---------------------------------------------------------------------------
\newcommand{\TelosProject}{Telos}
\newcommand{\TelosDocTitle}{}
\newcommand{\TelosDocSubtitle}{}
\newcommand{\TelosRevision}{}

\newcommand{\telosmeta}[4]{%
  \renewcommand{\TelosProject}{#1}%
  \renewcommand{\TelosDocTitle}{#2}%
  \renewcommand{\TelosDocSubtitle}{#3}%
  \renewcommand{\TelosRevision}{#4}%
  \hypersetup{pdftitle={#2},pdfsubject={#3; version #4},
    pdfkeywords={Telos, version #4},pdfauthor={Telos}}%
}

\pagestyle{fancy}
\fancyhf{}
\renewcommand{\headrulewidth}{0pt}
\renewcommand{\footrulewidth}{0.4pt}
\fancyfoot[L]{\footnotesize\textsc{\TelosProject}}
\fancyfoot[C]{\footnotesize\TelosDocTitle}
\fancyfoot[R]{\footnotesize\thepage}

% The first page of a one-pager carries the masthead instead of a running foot,
% so its footer is the compact provenance line.
\fancypagestyle{telosfirst}{%
  \fancyhf{}%
  \renewcommand{\headrulewidth}{0pt}%
  \renewcommand{\footrulewidth}{0.4pt}%
  \fancyfoot[L]{\footnotesize\textsc{\TelosProject}}%
  \fancyfoot[C]{\footnotesize\TelosRevision}%
  \fancyfoot[R]{\footnotesize\thepage}%
}

% ---------------------------------------------------------------------------
% Masthead. A heavy rule, the title, a light subtitle, a thin rule. The
% optional argument is a right-aligned tag (a species' scientific name, a
% lesson number) that sits on the title's baseline.
% ---------------------------------------------------------------------------
\newcommand{\telosmasthead}[1][]{%
  \thispagestyle{telosfirst}%
  \begingroup
  \setlength{\parskip}{0pt}%
  \noindent\rule{\linewidth}{1.6pt}\par
  \vspace{0.30em}
  \noindent
  \begin{minipage}[b]{0.66\linewidth}
    \raggedright{\LARGE\bfseries\TelosDocTitle}
  \end{minipage}%
  \hfill
  \begin{minipage}[b]{0.32\linewidth}
    \raggedleft{\small\itshape #1}
  \end{minipage}\par
  \vspace{0.18em}
  \noindent{\small\TelosDocSubtitle}\par
  \vspace{0.28em}
  \noindent\rule{\linewidth}{0.4pt}\par
  \endgroup
  \vspace{0.35em}
}

% ---------------------------------------------------------------------------
% Headings. Compact, small-caps, rule-anchored — a card has no room for the
% article class's default vertical air.
% ---------------------------------------------------------------------------
\titleformat{\section}{\normalfont\large\bfseries}{}{0pt}{}[\vspace{-0.55em}\rule{\linewidth}{0.4pt}]
\titlespacing*{\section}{0pt}{0.85em}{0.35em}
\titleformat{\subsection}{\normalfont\normalsize\bfseries}{}{0pt}{}
\titlespacing*{\subsection}{0pt}{0.6em}{0.2em}
\titleformat{\subsubsection}{\normalfont\small\bfseries\scshape}{}{0pt}{}
\titlespacing*{\subsubsection}{0pt}{0.45em}{0.12em}

% A short run-in heading for card interiors that must not consume a full line.
\newcommand{\telosrubric}[1]{\textbf{#1}\enspace}

% Cross-reference helper. Every Telos project is a set of loose printed sheets,
% so a reference names the sheet rather than a page number.
\newcommand{\sheetref}[1]{\textit{see the \textbf{#1} sheet}}

% Which decision, ADR or sheet governs the paragraph you are reading. A reader
% who disagrees with an instruction should be able to find the reasoning behind
% it rather than having to take it on trust.
\newcommand{\governedby}[1]{{\footnotesize\textsc{governed by #1}}}

% ---------------------------------------------------------------------------
% Quantities. siunitx is not in this TeX Live split, and the guides need only a
% fixed thin space and upright units, so define that directly rather than take
% a package dependency the documented install would not satisfy.
%   \qty{9}{kV}  ->  9 kV        \unitof{kV}  ->  kV
% ---------------------------------------------------------------------------
\newcommand{\unitof}[1]{\ensuremath{\mathrm{#1}}}
\newcommand{\qty}[2]{#1\,\unitof{#2}}
\newcommand{\numrange}[2]{#1\,--\,#2}
\newcommand{\qtyrange}[3]{#1\,--\,#2\,\unitof{#3}}
% Inches and feet appear constantly in the fishing guide; keep one spelling.
\newcommand{\inch}[1]{#1\,in}
\newcommand{\feet}[1]{#1\,ft}

% ---------------------------------------------------------------------------
% Cards and boxes.
% ---------------------------------------------------------------------------

% Titled card with a hairline frame. The default card for facts a reader looks
% up rather than reads through.
\newtcolorbox{teloscard}[1]{
  enhanced, breakable,
  colback=telospaper, colframe=telosink, colbacktitle=telospaper,
  coltitle=telosink,
  boxrule=0.5pt, sharp corners,
  left=6pt, right=6pt, top=4pt, bottom=4pt,
  toptitle=3pt, bottomtitle=3pt, titlerule=0.4pt,
  before skip=6pt, after skip=6pt,
  title={#1}, fonttitle=\small\bfseries
}

% Untitled left-ruled block: hierarchy without a redundant label.
\newtcolorbox{telosframe}{
  enhanced, breakable,
  colback=telospaper, frame hidden, boxrule=0pt, sharp corners,
  borderline west={1.3pt}{0pt}{telosink},
  left=8pt, right=2pt, top=3pt, bottom=3pt,
  before skip=6pt, after skip=6pt
}

% Filled emphasis panel — the "read this first" block. The 10% fill stays
% legible under a cheap laser toner curve and still reads as distinct.
\newtcolorbox{telospanel}[1]{
  enhanced, breakable,
  colback=telosfill, colframe=telosink, colbacktitle=telosfill,
  coltitle=telosink,
  boxrule=0.5pt, sharp corners,
  left=6pt, right=6pt, top=4pt, bottom=4pt,
  toptitle=3pt, bottomtitle=3pt, titlerule=0.4pt,
  before skip=6pt, after skip=6pt,
  title={#1}, fonttitle=\small\bfseries
}

% Hazard block. Deliberately the heaviest object on any page: a 1.4pt frame
% plus a solid title bar. Nothing else in Telos is allowed to look like this,
% so a reader skimming a lesson cannot mistake it for ordinary emphasis.
\newtcolorbox{teloshazard}[1]{
  enhanced, breakable,
  colback=telospaper, colframe=telosink,
  colbacktitle=telosink, coltitle=telospaper,
  boxrule=1.4pt, sharp corners,
  left=6pt, right=6pt, top=5pt, bottom=5pt,
  toptitle=3pt, bottomtitle=3pt,
  before skip=8pt, after skip=8pt,
  title={\MakeUppercase{#1}}, fonttitle=\small\bfseries
}

% A stop-and-check gate. Used before anything destructive, and before anything
% that would be expensive to discover later. Shared by every project that has
% a procedure someone could get hurt following carelessly.
\newtcolorbox{gate}[1]{
  enhanced, breakable,
  colback=telospaper, colframe=telosink,
  boxrule=1.0pt, sharp corners,
  left=6pt, right=6pt, top=4pt, bottom=4pt,
  toptitle=3pt, bottomtitle=3pt, titlerule=0.4pt,
  before skip=7pt, after skip=7pt,
  colbacktitle=telospaper, coltitle=telosink,
  title={\MakeUppercase{#1}}, fonttitle=\small\bfseries
}

% ---------------------------------------------------------------------------
% Rating dots. A five-step scale that prints identically in black-only: filled
% versus open circles, never a shaded bar.
% ---------------------------------------------------------------------------
\newcommand{\telosdot}{\tikz[baseline=-0.42ex]\fill[telosink] (0,0) circle (0.28ex);}
\newcommand{\telosopendot}{\tikz[baseline=-0.42ex]\draw[telosink,line width=0.28pt] (0,0) circle (0.28ex);}
\makeatletter
\newcounter{telos@rate}
\newcommand{\telosrating}[1]{%
  \begingroup
  \setcounter{telos@rate}{0}%
  \@whilenum\value{telos@rate}<#1\do{\telosdot\kern0.22em\stepcounter{telos@rate}}%
  \@whilenum\value{telos@rate}<5\do{\telosopendot\kern0.22em\stepcounter{telos@rate}}%
  \endgroup
}
\makeatother

% ---------------------------------------------------------------------------
% Tables.
%
% tabularx reads its body by scanning for a literal \end{tabularx}, so it can
% never be wrapped in a \newenvironment. Documents therefore open tabularx and
% longtable directly and take their house style from these column types and the
% booktabs rules. The one exception is the fact strip below, which is built on
% plain tabular precisely so it *can* be an environment.
% ---------------------------------------------------------------------------
\newcolumntype{L}{>{\raggedright\arraybackslash}X}
\newcolumntype{R}{>{\raggedleft\arraybackslash}X}
\newcolumntype{C}{>{\centering\arraybackslash}X}
\newcolumntype{B}{>{\bfseries\raggedright\arraybackslash}X}
\newcolumntype{N}{>{\raggedright\arraybackslash}p{0.16\linewidth}}

% Key/value fact strip. Two flush columns of short lookups; the label column is
% a fixed fraction so a stack of cards aligns across the whole guide.
\newenvironment{telosfacts}[1][0.24]
  {\begingroup\small
   \setlength{\parskip}{0pt}%
   \renewcommand{\arraystretch}{1.25}%
   \begin{tabular}{@{}>{\bfseries\raggedright\arraybackslash}p{\dimexpr#1\linewidth\relax}@{\hspace{0.7em}}>{\raggedright\arraybackslash}p{\dimexpr\linewidth-#1\linewidth-0.7em\relax}@{}}}
  {\end{tabular}\par\endgroup}
\newcommand{\telosfact}[2]{#1 & #2\\}

% ---------------------------------------------------------------------------
% Seasonal / diurnal activity bar.
%
% \telosseasonbar{4,3,2,1,1,2,3,3,4,5,4,3} draws a twelve-cell month strip
% where each value 0-5 is a fill height. Height, not gray value, carries the
% signal, so it is exact under any toner curve; the cell is outlined either way
% so an empty month is still visibly a month.
% ---------------------------------------------------------------------------
\newcommand{\telosseasonlabels}{J,F,M,A,M,J,J,A,S,O,N,D}
\newcommand{\telosseasonbar}[1]{%
  \begingroup
  \begin{tikzpicture}[x=1.25em,y=2.1ex,baseline=0pt]
    \foreach \v [count=\i from 0] in {#1} {
      \draw[telosrule,line width=0.3pt] (\i,0) rectangle (\i+0.86,5);
      \ifnum\v>0
        \fill[telosink] (\i,0) rectangle (\i+0.86,\v);
      \fi
    }
    \foreach \m [count=\i from 0] in \telosseasonlabels {
      \node[font=\tiny,anchor=north,inner sep=1.2pt] at (\i+0.43,-0.1) {\m};
    }
  \end{tikzpicture}%
  \endgroup
}

% Four-cell day-part strip: dawn, midday, dusk, night.
\newcommand{\telosdaylabels}{Dawn,Mid,Dusk,Night}
\newcommand{\telosdaybar}[1]{%
  \begingroup
  \begin{tikzpicture}[x=2.9em,y=2.1ex,baseline=0pt]
    \foreach \v [count=\i from 0] in {#1} {
      \draw[telosrule,line width=0.3pt] (\i,0) rectangle (\i+0.92,5);
      \ifnum\v>0
        \fill[telosink] (\i,0) rectangle (\i+0.92,\v);
      \fi
    }
    \foreach \m [count=\i from 0] in \telosdaylabels {
      \node[font=\tiny,anchor=north,inner sep=1.2pt] at (\i+0.46,-0.1) {\m};
    }
  \end{tikzpicture}%
  \endgroup
}

% A bar needs vertical room for its labels when it sits in a fact strip.
\newcommand{\telosbarrow}[1]{\rule{0pt}{4.2ex}#1\rule[-1.6ex]{0pt}{1.6ex}}

% ---------------------------------------------------------------------------
% Plates. A species or apparatus illustration with a caption rule beneath it.
% \telosplate{width fraction}{file}{caption}
% ---------------------------------------------------------------------------
\newcommand{\telosplate}[3]{%
  \begingroup
  \centering
  \includegraphics[width=#1\linewidth]{#2}\par
  \vspace{0.15em}
  \begin{minipage}{#1\linewidth}
    \rule{\linewidth}{0.4pt}\par
    \vspace{0.1em}
    {\footnotesize #3\par}
  \end{minipage}\par
  \endgroup
  \vspace{0.3em}
}

% TikZ drawing conventions shared by every rig, knot, and circuit diagram.
\tikzset{
  telos line/.style={draw=telosink,line width=0.7pt},
  telos hair/.style={draw=telosink,line width=0.35pt},
  telos heavy/.style={draw=telosink,line width=1.2pt},
  telos node/.style={draw=telosink,fill=telospaper,align=center,inner sep=4pt,font=\small},
  telos emphasis/.style={draw=telosink,line width=1.2pt,fill=telospaper,align=center,
    inner sep=4pt,font=\small\bfseries},
  telos label/.style={font=\scriptsize,inner sep=1.5pt},
  telos arrow/.style={-{Stealth[length=1.8mm]},draw=telosink,line width=0.6pt},
  telos dim/.style={|<->|,draw=telosink,line width=0.35pt},
  telos water/.style={draw=telosrule,line width=0.4pt,decorate,
    decoration={snake,amplitude=0.4mm,segment length=3mm}},
}

% ---------------------------------------------------------------------------
% Lesson furniture
% ---------------------------------------------------------------------------

% \lessonhead{number}{time}{who does what}
\newcommand{\lessonhead}[3]{%
  \par\noindent
  \begin{minipage}{\linewidth}
    \rule{\linewidth}{0.4pt}\par
    \vspace{0.3em}
    {\small\textbf{Lesson #1}\quad$\cdot$\quad #2\quad$\cdot$\quad #3\par}
    \vspace{0.25em}
    \rule{\linewidth}{0.4pt}
  \end{minipage}\par
  \vspace{0.5em}}

% The materials list. Two columns of short lines; a lesson that needs a long
% shopping list is a lesson that has not been thought through.
\newenvironment{materials}
  {\begin{telospanel}{What you need}\begin{multicols}{2}
   \begin{itemize}[leftmargin=1em,itemsep=0.05em,topsep=0pt]}
  {\end{itemize}\end{multicols}\end{telospanel}}

% The four rubrics every lesson carries, in the same order every time:
% what you are about to see, what to do, what it means, and what to ask.
\newcommand{\lessonsection}[1]{\section*{#1}\vspace{-0.35em}\rule{\linewidth}{0.4pt}\vspace{0.2em}}

% A prediction box. The point of the curriculum is that they commit to an
% answer before the demonstration, not after.
\newtcolorbox{predict}{
  enhanced, breakable,
  colback=telospaper, colframe=telosink,
  boxrule=0.5pt, sharp corners,
  left=6pt, right=6pt, top=4pt, bottom=4pt,
  before skip=6pt, after skip=6pt,
  borderline west={2.2pt}{0pt}{telosink},
  title={\small\bfseries Predict first}, fonttitle=\small\bfseries,
  colbacktitle=telospaper, coltitle=telosink,
  toptitle=3pt, bottomtitle=3pt, titlerule=0.4pt,
}

% ---------------------------------------------------------------------------
% Colophon. Rendered once, at the end of the last page, in the recovered footer
% area so it never creates a page of its own.
% ---------------------------------------------------------------------------
\newif\ifTelosColophonRendered
\newcommand{\TelosColophonExtra}{}
\newcommand{\TelosColophon}{%
  \ifTelosColophonRendered\else
  \global\TelosColophonRenderedtrue
  \par
  \enlargethispage{0.5in}%
  \vspace{0.3em}%
  \begingroup
  \setlength{\parskip}{0pt}%
  \begin{minipage}{\linewidth}
  \rule{\linewidth}{0.35pt}\par
  \vspace{0.3em}%
  \fontsize{7}{8.2}\selectfont
  \textbf{Telos.} \TelosProject{} — \TelosDocTitle. Revised \TelosRevision.
  \TelosColophonExtra{}
  Project-created text and diagrams are licensed \href{https://creativecommons.org/licenses/by/4.0/}{CC BY 4.0}.
  Incorporated public-domain artwork remains public domain and is credited on
  the plate it appears with. See \texttt{LICENSE} and \texttt{CREDITS.md} in the
  source. Nothing here is professional advice; verify current regulations and
  follow the stated safety procedure.
  \end{minipage}
  \endgroup
  \fi
}
\AtEndDocument{\TelosColophon}

% Shared macros for the homelab project.
%
% This project's documents are read in two very different situations: at a desk
% while designing, and standing in front of a broken machine at eleven at night
% with no network. The second case is the one that sets the style. Commands are
% set apart so they can be found by eye, every dangerous step has a visible
% verification after it, and nothing depends on being able to reach a web page.

% ---------------------------------------------------------------------------
% Decision-record entries, used by the generated digest.
% ---------------------------------------------------------------------------
\newcommand{\adrentry}[4]{%
  \par\vspace{0.7em}
  \needspace{4\baselineskip}%
  \noindent
  \begin{minipage}{\linewidth}
    \rule{\linewidth}{0.6pt}\par
    \vspace{0.25em}
    {\small\textbf{ADR #1}\quad #2\par}
    \vspace{0.1em}
    {\fontsize{7}{8.2}\selectfont\textsc{#3}\quad$\cdot$\quad #4\par}
    \vspace{0.2em}
    \rule{\linewidth}{0.3pt}
  \end{minipage}\par
  \vspace{0.3em}}

% ---------------------------------------------------------------------------
% Procedures.
%
% A rebuild manual is a sequence of commands and the observable evidence that
% each one worked. These environments keep that pairing structural rather than
% a matter of the author remembering.
% ---------------------------------------------------------------------------

% A block of shell to type. Deliberately not a listing package: plain verbatim
% in a ruled box prints identically everywhere and has no font dependency.
\newenvironment{shellblock}
  {\par\vspace{0.25em}\noindent
   \begin{minipage}{\linewidth}
   \rule{\linewidth}{0.3pt}\par\vspace{0.15em}
   \begingroup\footnotesize\ttfamily\obeylines\obeyspaces\parindent=0pt\parskip=0pt}
  {\endgroup\par\vspace{0.15em}\rule{\linewidth}{0.3pt}\end{minipage}\par\vspace{0.25em}}

% What you should see if the previous step worked. A step without one of these
% is a step you cannot verify, and in a rebuild manual that is a defect.
\newcommand{\expectthis}[1]{%
  \par\vspace{0.1em}
  \noindent\begin{minipage}{\linewidth}
  \hspace*{0.6em}\begin{minipage}{\dimexpr\linewidth-0.6em\relax}
  {\footnotesize\textbf{Expect:} #1\par}
  \end{minipage}\end{minipage}\par\vspace{0.3em}}

% \begin{gate} now lives in common/preamble.tex, shared with the other
% projects that document procedures.

% \governedby now lives in common/preamble.tex.

% An instance value that comes from the private overlay. In the published
% documents these render as visible placeholders, which is deliberate: a
% published manual must be complete and obviously generic, never look like it
% is describing somebody's actual network.
\newcommand{\instancevalue}[2]{\texttt{\textlangle #1\textrangle}%
  \footnote{Instance value. Supplied from \texttt{homelab/instance/} at build
  time; #2.}}
\newcommand{\placeholder}[1]{\texttt{\textlangle #1\textrangle}}

% ---------------------------------------------------------------------------
% Role and boundary diagrams.
% ---------------------------------------------------------------------------
\tikzset{
  role/.style={draw=telosink,line width=0.9pt,fill=telospaper,align=center,
    inner sep=5pt,font=\small,minimum width=2.6cm},
  roleoff/.style={draw=telosmute,line width=0.5pt,dash pattern=on 2.5pt off 1.8pt,
    fill=telospaper,align=center,inner sep=5pt,font=\small,minimum width=2.6cm},
  boundary/.style={draw=telosmute,line width=0.5pt,dash pattern=on 3pt off 2pt,
    rounded corners=3pt},
  flow/.style={-{Stealth[length=2mm]},draw=telosink,line width=0.7pt},
  hlabel/.style={font=\fontsize{6.4}{7.4}\selectfont},
}


\telosmeta{Homelab}{Provisioning Design}
  {Network boot, the authorization boundary, and how it is tested}{2026-07-26}

\begin{document}
\telosmasthead[High-level design]

A machine is installed by network booting it into an interactive environment,
answering questions at its console, and typing the target disk's serial number
to authorize the erase. There is no pre-staged record of what a machine should
become, no answer file, and no unattended mode. Physical access to the console
plus that typed confirmation \emph{is} the authorization model.

\section{The boot chain}
\begin{center}
\begin{tikzpicture}[x=1cm,y=1cm,
  pencil/.style={draw=telosink,line width=0.58pt,line cap=round,line join=round},
  faint/.style={draw=telosmute,line width=0.28pt,line cap=round},
  signal/.style={-{Stealth[length=1.8mm,width=1.2mm]},draw=telosink,
    line width=0.62pt,double=telospaper,double distance=0.65pt}]
  % A measured baseline and registration ticks make the plate read as a
  % technical field sketch rather than a row of generic application boxes.
  \draw[faint] (0,-0.93) -- (15.35,-0.93);
  \foreach \x in {0.25,3.95,7.15,10.25,13.45,15.10}
    \draw[faint] (\x,-0.86)--(\x,-1.00);

  % Target motherboard / UEFI ROM.
  \begin{scope}[shift={(0.25,0)}]
    \draw[pencil,rounded corners=1pt] (0,-0.58) rectangle (2.25,0.78);
    \draw[faint] (0.08,-0.49) rectangle (2.17,0.69);
    \draw[pencil] (0.28,-0.28) rectangle (1.10,0.43);
    \foreach \y in {-0.17,0.00,0.17,0.34}
      \draw[faint] (0.18,\y)--(0.28,\y) (1.10,\y)--(1.21,\y);
    \draw[pencil] (1.48,0.09) circle (0.29);
    \draw[faint] (1.27,0.09)--(1.69,0.09) (1.48,-0.12)--(1.48,0.30);
    \draw[pencil] (1.35,-0.40) rectangle (1.98,-0.23);
    \node[font=\fontsize{5.7}{6.2}\selectfont] at (0.69,0.08) {UEFI};
  \end{scope}

  % Bootstrap host: stacked network appliance, with the active DHCP service.
  \begin{scope}[shift={(3.95,0)}]
    \foreach \y in {-0.52,-0.10,0.32}
      {\draw[pencil,rounded corners=1pt] (0,\y) rectangle (2.15,\y+0.34);
       \draw[faint] (0.12,\y+0.11)--(1.38,\y+0.11);
       \fill (1.71,\y+0.17) circle (0.035);
       \draw[pencil] (1.86,\y+0.10) rectangle (2.02,\y+0.24);}
    \draw[faint] (0.20,-0.62)--(0.20,-0.52) (1.95,-0.62)--(1.95,-0.52);
  \end{scope}

  % The deliberately tiny TFTP payload: one ROM-sized loader.
  \begin{scope}[shift={(7.15,0)}]
    \draw[pencil] (0.18,-0.42) rectangle (1.72,0.58);
    \draw[faint] (0.28,-0.32) rectangle (1.62,0.48);
    \foreach \x in {0.36,0.65,0.94,1.23,1.52}
      \draw[pencil] (\x,-0.52)--(\x,-0.42) (\x,0.58)--(\x,0.68);
    \node[font=\fontsize{6.1}{6.4}\selectfont\ttfamily] at (0.95,0.16) {iPXE};
    \draw[faint] (0.46,-0.08)--(1.44,-0.08);
  \end{scope}

  % HTTP artifact and its manifest: the integrity relationship is visible.
  \begin{scope}[shift={(10.25,0)}]
    \draw[pencil] (0.08,-0.55)--(0.08,0.70)--(1.44,0.70)--(1.82,0.32)--(1.82,-0.55)--cycle;
    \draw[pencil] (1.44,0.70)--(1.44,0.32)--(1.82,0.32);
    \foreach \y in {0.13,-0.05,-0.23}
      \draw[faint] (0.32,\y)--(1.52,\y);
    \draw[pencil] (1.24,-0.38) circle (0.26);
    \draw[pencil] (1.12,-0.38)--(1.21,-0.29)--(1.38,-0.48);
  \end{scope}

  % The received installer environment, shown as a console on the target.
  \begin{scope}[shift={(13.45,0)}]
    \draw[pencil,rounded corners=1pt] (0,-0.43) rectangle (1.90,0.71);
    \draw[faint] (0.12,-0.31) rectangle (1.78,0.58);
    \draw[pencil] (0.28,0.35)--(0.47,0.20)--(0.28,0.05);
    \draw[pencil] (0.56,0.04)--(1.25,0.04);
    \draw[pencil] (0.69,-0.43)--(0.58,-0.68)--(1.32,-0.68)--(1.21,-0.43);
    \draw[pencil] (0.40,-0.68)--(1.50,-0.68);
  \end{scope}

  \foreach \a/\b in {2.50/3.95,6.10/7.33,8.87/10.33,12.07/13.45}
    \draw[signal] (\a,0.04)--(\b,0.04);

  \node[font=\small\bfseries,align=center] at (1.38,-1.22) {target firmware};
  \node[font=\small\bfseries,align=center] at (5.03,-1.22) {dnsmasq};
  \node[font=\small\bfseries,align=center] at (8.10,-1.22) {TFTP};
  \node[font=\small\bfseries,align=center] at (11.20,-1.22) {nginx};
  \node[font=\small\bfseries,align=center] at (14.40,-1.22) {installer};
  \node[hlabel,align=center] at (5.03,-1.63) {DHCP + PXE\\one authority};
  \node[hlabel,align=center] at (8.10,-1.63) {iPXE only\\unauthenticated};
  \node[hlabel,align=center] at (11.20,-1.63) {artifacts + manifest\\SHA-256 checked};
  \node[hlabel,align=center] at (14.40,-1.63) {Archiso\\interactive};
\end{tikzpicture}
\end{center}

\begin{telosfacts}[0.20]
\telosfact{Why the split}{TFTP is unauthenticated and slow, so it carries only the first-stage loader. Kernels, initramfs images and the root filesystem travel over HTTP, where they can be checksummed against a published manifest. \governedby{ADR 0044, ADR 0048}}
\telosfact{One DHCP authority}{Never two on one segment, at any stage, including transitions. Where external infrastructure already owns DHCP, the bootstrap host runs in ProxyDHCP mode and supplies boot options only. \governedby{ADR 0008, ADR 0052}}
\telosfact{No menu}{The iPXE script boots exactly one thing. A menu with a single entry is a timeout waiting to select the wrong option.}
\telosfact{Checksums, not signatures}{During Milestone A artifact integrity comes from a published SHA-256 manifest that the installer verifies before use. Authenticated boot is Milestone B. \governedby{ADR 0043}}
\end{telosfacts}

\section{The installer}
\begin{tabularx}{\linewidth}{@{}p{0.05\linewidth}p{0.26\linewidth}L@{}}
\toprule
& \textbf{Stage} & \textbf{What happens}\\
\midrule
1 & Collect & Disks, interfaces and firmware read from the running machine behind a substitutable seam, so the judgement logic is testable without hardware.\\
2 & Refuse & Legacy BIOS stops the run outright. So does a machine with no eligible disk, or a Controller with no wired interface. \governedby{ADR 0019, ADR 0050}\\
3 & Ask & Only the questions that apply. Validators see the answers already given, so a rule is reported at the prompt that broke it.\\
4 & Validate & The network plan is checked again as a whole before the summary. \governedby{ADR 0045}\\
5 & Summarise & Profile, hostname, firmware, interface, the entered and derived network plan, the exact disk, and the layout to be written. \governedby{ADR 0004}\\
6 & Authorize & The operator types the target disk's \textbf{serial number}. Not \texttt{yes}. \governedby{ADR 0058}\\
7 & Execute & Steps run in order and stop at the first failure. Nothing is rolled back. \governedby{ADR 0059}\\
8 & Record & A non-secret manifest is written to the new root and echoed to the console. \governedby{ADR 0060}\\
9 & Power off & Never reboot. Relocation while powered off is the DHCP-conflict boundary. \governedby{ADR 0010}\\
\bottomrule
\end{tabularx}

\section{What cannot be offered}
Some things are excluded from the choices rather than merely warned about,
because a warning is something an operator can read past at one in the morning.

\begin{telosfacts}[0.24]
\telosfact{Removable disks}{The installer boots from removable media. Offering the stick it is running from is the most destructive available mistake, so it never appears in the list --- with the reason printed, so the operator can see it was considered and rejected.}
\telosfact{Disks with no serial}{ADR 0058 confirms the erase by typing the serial. A disk that reports none cannot be confirmed, so it must not be offered. That rule fell out of the confirmation design rather than being invented separately.}
\telosfact{Disks under 64 GiB}{Not a computed minimum. A floor below which the operator has almost certainly selected the wrong device. \governedby{ADR 0051}}
\telosfact{Wireless interfaces}{Bare-metal network installation is wired, and a Controller serving DHCP over wireless is not a supported design.}
\telosfact{An unplugged wired interface}{\emph{Is} offered. The Controller is routinely provisioned on one network and relocated while powered off, so the managed interface is often disconnected during installation. Carrier is noted here and required at first boot. \governedby{ADR 0010, ADR 0009}}
\end{telosfacts}

\section{How it is tested}
The acceptance harness attaches to a pseudo-terminal and answers the genuine
prompts the way a person would, including typing the serial. There is no
unattended code path for it to use, so there is nothing to abuse on real
hardware, and the program under test is byte-identical to the one that runs on
metal.

\begin{telosfacts}[0.20]
\telosfact{One source of truth}{The installer renders a prompt registry; the harness reads the same registry to know what to answer. A reworded prompt cannot desynchronize them. \governedby{ADR 0058}}
\telosfact{The harness refuses to guess}{A prompt that appears and is not in the script is an error. An authorization prompt with no scripted confirmation is an error. Guessing is how a harness ends up confirming an erase nobody scripted.}
\telosfact{Same driver, two targets}{Locally it drives the installer as a subprocess. In the lab it drives a QEMU guest whose serial console is the process's stdio. The installer cannot tell the difference. \governedby{ADR 0056}}
\telosfact{The topology enforces the boundary}{The virtual lab has no user networking and no default NIC --- only a socket segment between guests. ADR 0011's no-routing rule is a property of the lab, not a setting in it.}
\telosfact{Fresh firmware every run}{Each guest gets its own OVMF variable store, so boot entries from one run cannot make a later run pass.}
\end{telosfacts}

\begin{telospanel}{A bug this found}
The rule that matters most is that the Controller's address must never sit
inside the DHCP pool --- dnsmasq must not be able to lease the address its own
DNS answers on. Asked in ADR 0045's order, that violation surfaced two prompts
late, while the operator was being asked for the pool's last address and told
that a different field was wrong. The questions are therefore asked subnet,
pool, Controller address: last, so the error lands on the prompt that caused it.
Driving the real interface found that; no unit test would have.
\end{telospanel}

\begin{gate}{Milestone A is disposable by construction}
Every installation during the functional proof is labelled
\texttt{development-proof} in its manifest, requires the LUKS2 passphrase at
every boot, and carries no production data. That label is a recorded state, not
a remembered one, so later tooling can refuse to treat such a machine as
production. \governedby{ADR 0043, ADR 0060}
\end{gate}
\end{document}
